\documentclass{article}
\usepackage{amsmath, amssymb, amsthm}
\usepackage{hyperref}
\usepackage{geometry}
\geometry{margin=1in}

\title{Erd\H{o}s Problem \#933}
\date{Last updated: 2025-08-31}
\author{Source: erdosproblems.com}

\begin{document}
\maketitle

\noindent\textbf{Status:} OPEN \quad \textbf{No prize}

\noindent\textbf{Tags:} number theory

\medskip

\noindent\textbf{Problem Statement:}

If $n(n+1)=2^k3^lm$, where $(m,6)=1$, then is it true that\[\limsup_{n\to \infty} \frac{2^k3^l}{n\log n}=\infty?\]




Mahler proved (a more general result that implies in particular) that\[2^k3^l&lt;n^{1+o(1)}.\]Erd\H{o}s \cite{Er76d} wrote 'it is easy to see' that for infinitely many $n$\[2^k3^l&gt;n\log n.\]Steinerberger has noted a simple proof of this fact follows from taking $n=2^{3^r}$ for any integer $r\geq 1$, when $k=3^r$ and $l=r+1$.




References

[Er76d] Erd\H{o}s, P., Problems and results on number theoretic properties of consecutive integers and related questions . Proceedings of the Fifth Manitoba Conference on Numerical Mathematics (Univ. Manitoba, Winnipeg, Man., 1975) (1976), 25-44.

\noindent\textbf{Related OEIS sequences:} possible


\bigskip
\noindent\small{Source: \url{https://www.erdosproblems.com/933}}
\end{document}
